% Test fixture : enonce.tex bfcours minimal pour valider atomize/monolithize.
% Couvre : titres [opt], titres {obl}, environnements imbriqués (MultiColonnes,
% tcbenumerate, center, crep), EXO contenant des Definition (qui NE doivent
% PAS être extraites car non top-level), commentaires `%`, et un \input
% existant à préserver. Inclut aussi des EXO précédés de \def\rdifficulty
% pour valider l'absorption arrière.

\section{Premier chapitre}

\subsection{Notions de base}

\begin{Definition}[Variable aléatoire réelle]
Une\voc{variable aléatoire réelle} associe à chaque issue un \acc{nombre réel}.
\begin{MultiColonnes}{2}
    \tcbitem $X : \omega \mapsto X(\omega) \in \R$
    \tcbitem On dit \acc{discrète} si finie.
\end{MultiColonnes}
\end{Definition}

\begin{Propriete}{Existence et unicité}
Il existe une \emph{unique} fonction $f$ telle que $f' = f$ et $f(0) = 1$.
\end{Propriete}

% Le commentaire suivant ne doit PAS perturber le parser :
% \begin{Definition}[Fausse définition commentée] ... \end{Definition}

\begin{Exemple}[Pièce truquée -- deux lancers]
\begin{tcbenumerate}[2]
    \tcbitem $P(X = 0) = \dfrac{4}{9}$.
    \tcbitem $P(X = 1) = \dfrac{4}{9}$.
\end{tcbenumerate}
\vspace{-0.3cm}\end{Exemple}

\subsection{Exercices}

\def\rdifficulty{1}
\begin{EXO}{Calcul d'espérance}{1SP23}
\begin{Definition}[Définition imbriquée dans EXO]
Cette définition est dans un EXO : elle NE doit PAS être atomisée séparément.
Elle doit toutefois suivre l'EXO dans son fichier atomisé.
\end{Definition}

\begin{tcbenumerate}[2]
    \tcbitem \tcbitempoint{2}\acc{Calculer} $E(X)$.
    \tcbitem \tcbitempoint{1}\acc{Conclure}.
\end{tcbenumerate}

\exocorrection

$E(X) = 2$.
\end{EXO}

\def\rdifficulty{2.5}
\begin{EXO}{Tirage compliqué}{1SP24}
Un EXO avec sa difficulté juste au-dessus.
\exocorrection
La correction.
\end{EXO}

\section{Deuxième chapitre}

\input{sections/Definition/def_existante.tex}

\begin{Remarque}
Une remarque sans titre.
\end{Remarque}

\begin{EXO}{EXO sans rdifficulty}{}
Cet EXO n'a pas de \verb|\def\rdifficulty| au-dessus.
\end{EXO}

\begin{Demonstration}[Cas n°2]
On a $\Phi(x) = 0$ donc \encadrer[defi]{$f$ ne s'annule pas}.
\end{Demonstration}
